% Generated draft: requires mathematical and source review.
\documentclass[11pt]{article}
\usepackage[margin=25mm]{geometry}
\usepackage{fontspec}
\setmainfont{Latin Modern Roman}
\usepackage{amsmath,amssymb,mathtools,unicode-math}
\setmathfont{Latin Modern Math}
\usepackage{enumitem,needspace,xurl,hyperref}
\setlength{\parindent}{0pt}
\setlength{\parskip}{5pt}
\setlength{\emergencystretch}{3em}
\newcommand{\N}{\mathbb{N}}
\newcommand{\Z}{\mathbb{Z}}
\newcommand{\Q}{\mathbb{Q}}
\newcommand{\R}{\mathbb{R}}
\newcommand{\C}{\mathbb{C}}
\newcommand{\F}{\mathbb{F}}
\newcommand{\A}{\mathbb{A}}
\newcommand{\PP}{\mathbb{P}}
\newcommand{\E}{\mathbb{E}}
\newcommand{\eps}{\varepsilon}
\newcommand{\dd}{\,\mathrm{d}}
\begin{document}
\section*{UnsolvedMath ID: 30001280}
\textbf{Source problem code:} catalogue entry
\subsection{Definitions and mathematical statement}
For a unital ring \(A\), write \(A^\times\) for its group of units and define its algebraic \(K_1\)-group by
\[
K_1(A)=\mathrm{GL}_\infty(A)^{\mathrm{ab}},
\qquad
\mathrm{GL}_\infty(A)=\varinjlim_n\mathrm{GL}_n(A),
\]
where stabilization sends \(M\) to \(\operatorname{diag}(M,1)\). For a profinite group \(H\) and a unital coefficient ring \(O\), the usual inverse-limit construction is
\[
\Lambda_O(H)=\varprojlim_{U\trianglelefteq_{\mathrm{open}}H}O[H/U],
\]
with transition maps induced by quotient homomorphisms. The supplied record does not specify the coefficient ring, its topology, or the hypotheses on the groups needed for the intended Iwasawa-theoretic setting. For a profinite group \(H\), its topological abelianization is \(H^{\mathrm{ab}}=H/\overline{[H,H]}\); whether this is the intended convention must be confirmed from the source.

The supplied computational request is to determine the relevant \(K_1\)-groups and characterize the image of a map displayed as
\[
\theta:K_1\bigl(\Lambda_O(G)\bigr)
\longrightarrow
\prod_{i\geq 0}\Lambda_O\bigl(G_i^{\mathrm{ab}}\bigr)^\times,
\]
and to give the corresponding characterization for a localized map denoted \(\theta_S\). Here the index range \(i\geq 0\) is retained exactly as supplied. The record does not define \(O\), \(G\), the family \((G_i)_{i\geq 0}\), the construction of \(\theta\), the localization set \(S\), or the domain and codomain of \(\theta_S\). It also does not identify the full collection of ``relevant Iwasawa algebras.'' Consequently, the displayed request is a schematic research-program description, and no definite image conditions, quantified conjecture, or complete standalone mathematical problem can be recovered from this record.

The supplied verification note states that the report presents the computation and theta-image characterization as Theorem 1.1 and its localized consequence, rather than as an open problem. This is a supplied source-verification claim, not an independently checked theorem statement. The registry nevertheless retains the status \emph{open}, marks the statement \emph{unrecoverable}, and supplies no clean replacement statement. These qualifications remain unresolved.
\subsection{Short English statement}
The record asks for computations of \(K_1\)-groups of Iwasawa algebras and descriptions of the images of \(\theta\) and its localized version \(\theta_S\). Essential objects, hypotheses, and maps are unspecified. Moreover, the supplied verification note identifies the material as a theorem and a localized consequence, while the registry retains an open label. A precise standalone problem therefore cannot be stated from the supplied evidence.
\subsection{Sources}
\begin{itemize}
\item Supplied local UnsolvedMath registry record, permanent ID \(30001280\), exact source problem code \texttt{catalogue entry}, titled \(K_1\) Groups and Theta Images of Iwasawa Algebras; collection: Oberwolfach Reports Open Problems. The record lists the year 2009 and supplies no named proposer.
\item \emph{Algebraische Zahlentheorie} (2009), Oberwolfach Reports, pp.\ 1667--1730, DOI \texttt{10.4171/owr/2009/30}. Supplied source URL and literature-source link, labeled ``OWR report'': \url{https://doi.org/10.4171/owr/2009/30}. The supplied background attributes the extraction to workshop reports and their open-problem discussions; no problem-specific page or author is provided.
\item Supplied literature assessment, checked on \texttt{2026-08-21}: the status is labeled open, but the extraction is described as a program of noncommutative Iwasawa computations rather than a precise resolved claim, and no standalone complete problem is certified solved. This is supplied dated curation, not independently verified; the supplied background calls for independent verification of time-sensitive status and cited claims.
\item Supplied statement-verification note: the official report is said to present the computation and theta-image characterization as Theorem 1.1 and its localized consequence, not as an open problem. The statement is marked unrecoverable, no clean statement is promoted, and the original record is retained for provenance. The report has not been independently consulted for this draft.
\end{itemize}
\end{document}
